% ======================= 结 语 =======================
\clearpage
\chapter*{结\quad 语}
\phantomsection
\addcontentsline{toc}{chapter}{结\quad 语}
\markboth{结\quad 语}{结\quad 语}

书到这里算是写完了。回到东京以后，我将从莫斯科带回来的那半本笔记重新誊过一遍，也将早年自己画过、写过的一摞月亮旧稿收进了纸箱底下。那些稿子当年在读者圈里流传得很广，大家读着开心，我也以为只是自己随手编出的一套规矩；直到对照着几百片合金板和湖底孢粉表逐行往下核，我才弄明白，真正撑住那座城市的，只是一份两万年都没有人肯去修改的旧规程，与浪漫的幻想或高深的奇迹都扯不上关系。

在著者中文版序言里，我曾提前交代过一句：全书有一处推论超出了手头材料所能承担的分量，相关限度须留到结语里交待清楚。这一处，便落在第二章末尾。在那一节里，我把贝加尔湖六点零三米深的孢粉断层、题头上的「计口三十一万二千」、有穷氏地下的金属空腔，还有月球的转动惯量与平均密度摆在一张桌面上，紧接着便落笔写下「月球更像是一艘跑完三十八万公里航程、停泊在现今轨道上的船」。

现在把话摊开来说，这一步确实迈得过头了。我没去过月面，手里也没有钻穿过那层不知究竟有多厚的月壳；档案里找不到半张船坞清单、半页推进机参数，或者任何一帧记录它如何离开地面轨道的航迹图。把现有材料老老实实摆成一排，硬据能支撑的底部只到「地面距今两万一千三百年确曾有一次极快的植被剧变」「城中留下的居止条文写明居住者身处壳内」「月壳之下留有大尺度的空腔与人工环制」这三层为止。再往外推一步，断定那颗天体本身就是从贝加尔湖畔起航、跑完三十八万公里停进今天这条轨道的载具，更多是为了把两头零散的材料接上一条连贯的线索，超出了手头物证所能承担的分量。后面各章之所以能接连往下走，很大程度上也是因为暂时借用了这一步推论当作脚手架。读者若只认物证、不认这半步跳跃，第二章末尾那句话完全可以留待日后出现更硬的深层钻探记录再来改写。

即便把这一层推论的水分挤干，十二章收尾时依然留着一件确凿的事实：自公元前二二一年生态崩溃、六万五千人仓促撤离之后，月球上再也没有别的月之民上去过。那座留在正面一角的静海七号跟背面的六道居住环一道熄灯封存，唯一留在这颗星球上的，只剩公元前一九八五年从地面接回、随后在断电隔离舱里无人过问的那一个服药者，以及两万年来没人敢去动过一笔的空置规程。

去年秋天我回了一次长野老家。傍晚在山边那座旧神社的石阶上站了一会儿，抬头刚好能看见东边升起来的月亮。从神社所见之月，照旧是那么亮、那么安静，跟五十年前我还是孩子时在同一个院子里见过的并无两样。不过现在站在阶下抬头往上看，脑子里想的已经再难回到旧日那些故事里去了：我知道那张始终朝向这一侧的月面背后，曾经画过几圈规整得反常的低地，而正面某块荒凉的月海边缘，至今还横着几道两万年前按秦汉尺度量出来的空墙。

我只是一个照着材料工作的人，既非历史学者，也非天体物理学家；材料摆在桌上长什么样，我只照着尺寸量出来、照着字口抄下来。写完这十二章，桌面上依然留着八千个找不着去向的人口、一千个对不齐的登船尾数，还有一连串写着「未解」与「不答」的空格。把这些空白原封不动地摆在纸面上，比替它们编一个圆满收尾更接近真实。

\begin{flushright}
{\small 太田顺也}\\[2mm]
{\footnotesize\color{pencil} 二〇二四年十二月，于东京}
\end{flushright}

% ======================= 附录一 =======================
\clearpage
\chapter*{附录一\quad 本书大事年表}
\phantomsection
\addcontentsline{toc}{chapter}{附录一\quad 本书大事年表}
\markboth{附录一　本书大事年表}{附录一　本书大事年表}

本表依全书各章考订次序排列，用以贯通月面与地面两端纪年。凡传世文献异说或城内纪年无法直接折算公元者，各据正文章节与《月之都——〈广寒宫残简〉辑佚与考释》附表注明依据。

\vspace{2mm}
{\footnotesize\renewcommand{\arraystretch}{1.20}
\begin{longtable}{@{}p{32mm}p{60mm}p{40mm}@{}}
\toprule
年代 & 事件 & 主要依据 \\
\midrule
\endfirsthead
\toprule
年代 & 事件 & 主要依据 \\
\midrule
\endhead
\bottomrule
\endfoot
距今约 21{,}300 年 & 地面生态与气候突变，伴随技术退步及持续减员 & 贝加尔湖岩芯 \texttt{ФВ-1957-04}（6.03 м 层）及 BDP-93-2 \\
其后数百年 & 月船建造与迁入；开工计口 312{,}000，完工登船 150{,}000 & 《章程总目》题头（\texttt{ФВ-1965-09}） \\
约前 2070 年 & 夏立国；月之都启用登船前所立社会实验判据 & 《史记·夏本纪》；规程第四十四枚 \\
约前 2050 年 & 嫦娥下地，任「地之主任」兼「手之总管」 & 外派名册（第四章）；净秽机制 \\
约前 2000 年 & 删三术劣化历法下送中原；后羿代夏（十日射九日）、改历 & 《左传》襄公四年；规程第二、三、二十七枚 \\
约前 1995 年 & 后羿历五年自行补全三术，重现「天文密葬法」隐层 & 规程第二十七枚背面 \\
约前 1990 年 & 上行「底牌见于人手」；嫦娥经信道同时下发处置令，后羿与伯封互杀同死 & 规程第四十枚；《左传》昭公二十八年 \\
约前 1990 年 & 纯狐借旧历记号与地下设备执行纯化，未退出而同死；档注「未解」 & \texttt{ФВ-1972-06}、\texttt{ФВ-1972-07} \\
约前 1985 年 & 嫦娥被接回月面，封入独立隔离舱段，始记「创面不合」 & \texttt{ФВ-1969-03}；《淮南子》 \\
约前 1980 年 & 月兔之乱，月兔全灭，月之民死 22{,}000；逆向还原净秽山库存（「出之者，封」）；山内改积原状污秽 & 《广寒宫残简》简十九；\texttt{ФВ-1986-04} 旁证 \\
约前 1900 年 & 少康复国（地面土著自行完成中兴） & 《左传》哀公元年 \\
约前 455 年 & 污秽近临界、地热衰减初显；第一批约 6{,}000 人降落 & \texttt{ФВ-1970-18}；《广寒宫残简》简二十九 \\
约前 405—前 255 年 & 第二至第五批依计划降落（合计约 49{,}000 人） & \texttt{ФВ-1970-18}；《广寒宫残简》简三十 \\
前 221 年 & 末期约二百年地热累计降 3.1\%，生态突发崩溃；第六批 65{,}000 人撤离（实载合 120{,}000、在册 119{,}000）；玉山（4{,}713 м）与十二舱落 34.6°N、110.4°E；隔离舱记录至「酷似蟾蜍之物」而断；「绝地天通」截断垂直通道 & \texttt{Ф.\,117, оп.\,4, д.\,2187}、\texttt{ФВ-1969-03}、\texttt{ФВ-1970-18}、\texttt{ФВ-1971-22}；《广寒宫残简》简二十至三十一 \\
前 211 年 & 东郡「坠星」至地为石（晚期着陆残片） & 《史记·秦始皇本纪》；\texttt{ФВ-1974-12} \\
前 210 年 & 徐福东渡；八意永琳、蓬莱山辉夜等 217 名月之民随船往日本 & 《史记·淮南衡山列传》；\texttt{ФВ-1978-11} \\
712／720 年 & 《古事记》《日本书纪》成书；东渡口传凝固为「天津神」 & 记纪 \\
九至十世纪 & 《竹取物语》成书 & 传世本 \\
1959—1986 年 & 轨道影像、合金册片、有穷氏故地、二里头柱洞及净秽山钻探先后编目 & \texttt{ФВ-1963-07} 至 \texttt{ФВ-1986-04} \\
\end{longtable}}

% ======================= 附录二 =======================
\clearpage
\chapter*{附录二\quad 本书所用档案编号清单}
\phantomsection
\addcontentsline{toc}{chapter}{附录二\quad 本书所用档案编号清单}
\markboth{附录二　档案编号清单}{附录二　档案编号清单}

本书正文及注释征引的基金会与相关机构档案，依编号顺序汇列如下；第 117 全宗档案与各地出土卷宗互校使用，一并附列于末。

\vspace{2mm}
{\footnotesize\renewcommand{\arraystretch}{1.20}
\begin{longtable}{@{}p{34mm}p{52mm}p{36mm}p{12mm}@{}}
\toprule
档案编号 & 内容 & 藏所／整理单位 & 出现处 \\
\midrule
\endfirsthead
\toprule
档案编号 & 内容 & 藏所／整理单位 & 出现处 \\
\midrule
\endhead
\bottomrule
\endfoot
\texttt{ФВ-1957-04} & 贝加尔湖高沉积率外缘活塞岩芯及六点零三米微采复核档 & VFS 基金会异常文献部 & 第二章 \\
\texttt{ФВ-1963-07} & 远东委员会异常金属样品分析卷宗（含鄂尔多斯 \textsc{э}-14） & VFS 基金会异常文献部 & 第十一章 \\
\texttt{ФВ-1963-17} & 静海七号移交清单与测量附表 & VFS 基金会异常文献部 & 第一章 \\
\texttt{ФВ-1963-46} & 静海西缘铼锇同位素异常金属圆片标本 & VFS 基金会异常文献部 & 第一章 \\
\texttt{ФВ-1965-09} & 《章程总目》（四百一十二枚，六卷） & VFS 基金会异常文献部 & 第二至五章 \\
\texttt{ФВ-1968-14} & 《工籍残卷》（三枚：轮值、工别、呈报） & VFS 基金会异常文献部 & 第四、十章 \\
\texttt{ФВ-1968-29} & 深地热流与轨道耦合遥测备忘 & VFS 基金会异常文献部 & 第十一章 \\
\texttt{ФВ-1969-03} & 《隔离舱管理档》（十四块薄合金板） & VFS 基金会异常文献部 & 第九章 \\
\texttt{ФВ-1970-18} & 《撤离签符与登船核对册》（附青铜签符） & VFS 基金会异常文献部 & 第十一、十二章 \\
\texttt{ФВ-1971-08} & 《规程残卷》（十一枚：遣物、信道、密葬、处置） & VFS 基金会异常文献部 & 第七、八章 \\
\texttt{ФВ-1971-22} & 河南偃师二里头遗址三区 T-14 九柱洞群与朱书陶片档案 & VFS 基金会异常文献部 & 第十一章 \\
\texttt{ФВ-1972-06} & 有穷氏故地地下空腔测绘档案 & VFS 基金会异常文献部 & 第二、八章 \\
\texttt{ФВ-1972-07} & 有穷氏旧历竹片抄本（三十一片，存十八片可辨） & VFS 基金会异常文献部 & 第八章 \\
\texttt{ФВ-1974-12} & 东郡旧河床熔蚀合金残片检验报告 & VFS 基金会异常文献部 & 第十一章 \\
\texttt{ФВ-1978-11} & 北九州弥生中期湿地木简（记二百一十七口及永琳、辉夜、月夜见） & VFS 基金会异常文献部 & 第十二章 \\
\texttt{ФВ-1979-08} & 静海西缘近地遥测校对叠图（\textsc{mt}-07 八边形异常区） & VFS 基金会异常文献部 & 第一章 \\
\texttt{ФВ-1986-04} & 净秽山北侧浅层钻孔岩芯与四十六枚断钻头记录 & VFS 基金会异常文献部 & 第六、十章 \\
\texttt{ФВ-1989-12} & 地球同步轨道未编目钛铝钒碎片回收与坑蚀定年档案 & VFS 基金会异常文献部 & 第一章 \\
\texttt{Ф.\,117, оп.\,4, д.\,2187} & 静海七号原始图纸、退役清单、玉山再入及认同留存率统计转抄表 & 俄罗斯科学院档案馆 & 第一、十一、十二章 \\
\end{longtable}}

% ======================= 附录三 =======================
\clearpage
\chapter*{附录三\quad 术语对照表（中—俄—英）}
\phantomsection
\addcontentsline{toc}{chapter}{附录三\quad 术语对照表（中—俄—英）}
\markboth{附录三　术语对照表}{附录三　术语对照表}

本表收录本书所用核心概念与专名。与《月之都——〈广寒宫残简〉辑佚与考释》及丛书体例共用的术语，俄、英对译一律照丛书定名；「说明」栏只注明译名取舍理由，不重述定义。

\vspace{2mm}
{\footnotesize\renewcommand{\arraystretch}{1.12}
\begin{longtable}{@{}p{25mm}p{35mm}p{35mm}p{37mm}@{}}
\toprule
中译 & 俄文 & 英文 & 说明 \\
\midrule
\endfirsthead
\toprule
中译 & 俄文 & 英文 & 说明 \\
\midrule
\endhead
\bottomrule
\endfoot
月之都 & Лунная столица & Lunar Capital & 兼指背面本体与正面静海七号 \\
月之民 & лунный народ & lunar folk & 废旧称「月面的居民」，照丛书统一定名 \\
华夏神族 & божественный род Китая & the Divine Lineage & 仅指中原十二万降落者及其后裔 \\
天津神 & амацуками & amatsukami & 保留日文音译，专指与东渡支系合流之名号 \\
污秽 & скверна & impurity & 取「变化、流转」之义，不用现代卫生词汇 \\
净秽山 & гора Очищения & Mount of Purification & 照汉文原名意译 \\
蓬莱药 & эликсир Хорай & Hourai elixir & 沿用日文原著汉字定译 \\
蓬莱人 & хорайский бессмертный & Hourai immortal & 与档案「服药者」对应，区分一般月之民 \\
月兔 & лунный кролик & moon rabbit & 本书正名；残简考释本只取「捣者」 \\
手之总管 & смотритель «рук» & overseer of Hands & 照《名籍》卷原称直译 \\
地之主任 & земной уполномоченный & terrestrial director & 统一全书官称，不混作地面负责人 \\
纯化 & очистка-разложение & purification-dissolution & 兼顾分解与剥离原整体两层含义 \\
天文密葬法 & метод сокрытия светил & celestial concealment & 依「密葬」隐没星位本义意译 \\
昆仑山号 & «Куньлунь» & Kunlun Vessel & 撤离所用留轨母船名称 \\
玉山 & Нефритовая гора & Jade Mountain & 沿用着陆载具汉文旧称 \\
绝地天通 & пресечение небесного пути & severing of the skyway & 依《国语》旧名，指截断垂直运输通道 \\
记忆消除场 & поле забвения & field of oblivion & 不译「遗忘场」，从丛书统一定名 \\
生态维持系统 & система поддержания экоса & ecosphere maintenance system & 照丛书工程类术语通例 \\
\end{longtable}}

% ======================= 参考文献 =======================
\clearpage
\chapter*{参\quad 考\quad 文\quad 献}
\phantomsection
\addcontentsline{toc}{chapter}{参\quad 考\quad 文\quad 献}
\markboth{参考文献}{参考文献}

\begingroup
\footnotesize
\linespread{1.25}\selectfont
\setlength{\parindent}{0pt}\setlength{\parskip}{1.6pt}

一、中日传世文献与点校本
\par\vspace{1.5mm}
\begin{list}{}{\leftmargin=1.8em \itemindent=-1.8em \topsep=0pt \parsep=1.4pt \itemsep=0pt}
\item 《尚书》（《十三经注疏》本），中华书局 1980 年影印本。
\item 《诗经》（《十三经注疏》本），中华书局 1980 年影印本。
\item 徐元诰《国语集解》，中华书局 2002 年点校本。
\item 《春秋左传正义》（《十三经注疏》本），中华书局 1980 年影印本。
\item 洪兴祖《楚辞补注》，中华书局 1983 年点校本。
\item 郭庆藩《庄子集释》，中华书局 1961 年点校本。
\item 王先谦《荀子集解》，中华书局 1988 年点校本。
\item 刘文典《淮南鸿烈集解》，中华书局 1989 年点校本。
\item 袁珂《山海经校注》，上海古籍出版社 1980 年版。
\item 《史记》，中华书局 1959 年点校本。
\item 《汉书》，中华书局 1962 年点校本。
\item 《后汉书》，中华书局 1965 年点校本。
\item 《古事记》（《日本古典文学大系》本），岩波书店 1958 年版。
\item 《日本书纪》（《日本古典文学大系》本），岩波书店 1967 年版。
\item 《竹取物语》（《新日本古典文学大系》本），岩波书店 1997 年版。
\end{list}

\vspace{4mm}
二、相关专著及专题文献
\par\vspace{1.5mm}
\begin{list}{}{\leftmargin=1.8em \itemindent=-1.8em \topsep=0pt \parsep=1.4pt \itemsep=0pt}
\item 巴兰尼科夫辑校《月之都——〈广寒宫残简〉辑佚与考释》，李文亚、张建国译，莫斯科大学出版社 2024 年版。
\item 索科洛夫《记忆消除场通论》，莫斯科大学出版社。
\item 中国科学院考古研究所二里头工作队：《河南偃师二里头遗址三、八区发掘简报》，《考古》1975 年第 5 期（作地层背景参证；异常柱洞及金属片见 \texttt{ФВ-1971-22}）。
\item Баранников А. В. \textit{Атлас реконструкций лунных объектов}. СПб., 2016.
\item Баранников А. В. \textit{Лунная прародина: опыт реконструкции}. М.: Изд-во МГУ, 2018.
\item Соколов Б. П. \textit{Осадочные керны Байкала}. М.: Наука, 1961.
\item Prokopenko, A. A., Karabanov, E. B., Williams, D. F. et al. \textit{The detailed record of climatic events during the past 75,000 yrs BP from the Lake Baikal drill core BDP-93-2}. Quaternary International 80–81 (2001): 59–68.
\item ESA. \textit{ESA’s Annual Space Environment Report}. Darmstadt, 2019（仅参证地球同步轨道未编目碎片背景）.
\item Wilhelms, D. E. and McCauley, J. F. \textit{Geologic Map of the Near Side of the Moon}, USGS Map I-703. Washington, D.C., 1971（\textsc{mt}-07 异常轮廓出内部叠图 \texttt{ФВ-1979-08}）.
\end{list}
\endgroup

% ======================= 索 引 =======================
\clearpage
\chapter*{索\quad 引}
\phantomsection
\addcontentsline{toc}{chapter}{索\quad 引}
\markboth{索\quad 引}{索\quad 引}

{\footnotesize\color{pencil}（页码为本书阿拉伯数字正文页码。）\par{}}
\vspace{2mm}

\begingroup
\footnotesize\linespread{1.35}\selectfont
\newcommand{\idx}[2]{\hangindent=1.4em \noindent #1\dotfill #2\par{}}
\begin{multicols}{2}
\idx{八意永琳}{78—79、82、87}
\idx{半坡}{1、72}
\idx{贝加尔湖}{1—3、11、13、17、39、83—84、86}
\idx{伯封}{48、53、86}
\idx{嫦娥}{28、43、47、49、52—53、57—59、62、79—80}
\idx{纯化}{2、47、49、51—53、55—56、86、90}
\idx{纯狐}{2、48—50、52—54、56—57、86}
\idx{次生黄土}{2—3、71、73}
\idx{地核}{60、68、71}
\idx{地之主任}{28、43、49、54、62、80、86、90}
\idx{东渡}{2、61、64、75、77—80、87、90}
\idx{东郡坠星}{77}
\idx{二里头遗址}{72、88、92}
\idx{隔离舱}{1—2、52、54—60、62、69、85—88}
\idx{工籍残卷}{1、25—28、63、65、88}
\idx{规程残卷}{1—2、40—41、46、50—51、88}
\idx{华夏神族}{2、75、80—82、90}
\idx{后羿}{2、14、44—46、48—49、53、59、78、86}
\idx{记忆消除场}{18、39、58、74、79、91}
\idx{净秽卷}{29—31}
\idx{净秽山}{37、62、64、67}
\idx{静海七号}{7—9、20、62}
\idx{居止卷}{17}
\idx{绝地天通}{2、67、72—73、91}
\idx{昆仑山号}{69、91}
\idx{六批}{2、64、67、73、87}
\idx{名籍卷}{14、26、69}
\idx{蓬莱人}{57、90}
\idx{蓬莱山辉夜}{78、80、82、87}
\idx{蓬莱药}{36—39、57—59、62、64、78、90}
\idx{岁计卷}{30}
\idx{天外认同留存率}{2、76—77}
\idx{天津神}{79—80、90}
\idx{天文密葬法}{45、61、78、86、90}
\idx{西王母}{54、59、61、79}
\idx{徐福}{2、37、77、82、87}
\idx{月读命}{79—80}
\idx{月兔}{2、36、38—39、47、59—60、62—66、79、86}
\idx{月兔之乱}{2、38—39、62—63、65、79、86}
\idx{月夜见}{78—79}
\idx{月之民}{20、27、63、69、77—79}
\idx{有穷氏}{1—2、14—15、20、44、47、49—51、84、87—88}
\idx{玉山}{2、70—71、73、87、89、91}
\idx{章程总目}{13、25—27、36、40、63、68}
\end{multicols}
\endgroup

\vspace{8mm}
\begin{center}
{\small\color{pencil}——\quad 全书终\quad ——}
\end{center}
